\documentclass[letterpaper,11pt]{article}

\usepackage[utf8]{inputenc}
\usepackage[T1]{fontenc}
\usepackage{lmodern}
\usepackage[spanish,es-tabla,es-nodecimaldot]{babel}
\usepackage[margin=2.5cm]{geometry}
\usepackage{setspace}
\setstretch{1.15}
\usepackage{graphicx}
\usepackage{booktabs}
\usepackage{tabularx}
\usepackage{amsmath}
\usepackage{float}
\usepackage[font=small,labelfont=bf]{caption}
\usepackage{csquotes}
\usepackage[hidelinks]{hyperref}

% Estilo de referencias: APA 7
\usepackage[style=apa,backend=biber]{biblatex}
\addbibresource{referencias.bib}
\renewcommand*{\bibfont}{\small}

% --- Complementos del resumen ejecutivo ---
\usepackage{xcolor}
\setlength{\parindent}{1.5em}

% --- Encabezado: logo del Tec solo en la primera página ---
\usepackage{fancyhdr}
\geometry{headheight=52pt,headsep=0.3cm}
\renewcommand{\headrulewidth}{0pt}
\newcommand{\logo}[2]{%
  \IfFileExists{#1}%
    {\includegraphics[height=1.2cm]{#1}}%
    {\fbox{\parbox[c][1.1cm][c]{3cm}{\centering\sffamily\small #2}}}}
\fancypagestyle{primera}{%
  \fancyhf{}%
  \fancyhead[L]{\logo{logo.jpeg}{Logo Tec}}%
  \fancyfoot[C]{\thepage}}
\pagestyle{fancy}
\fancyhf{}
\fancyfoot[C]{\thepage}

% --- Recuadros de cifras (texto grande) ---
% \cifra{fondo}{color del texto}{valor}{etiqueta}{unidad}
\newcommand{\cifra}[5]{%
  \begingroup\setlength{\fboxsep}{5pt}%
  \colorbox{#1}{\parbox[t]{\dimexpr\linewidth-2\fboxsep\relax}{%
    \setstretch{1}\color{#2}{\Large\bfseries #3}\par
    #4%
    \ifx\relax#5\relax\else\par{\footnotesize #5}\fi}}%
  \endgroup}
\newcommand{\cifragris}[4]{\cifra{black!8}{black}{#1}{#2}{#3}}
\newcommand{\cifraacento}[4]{\cifra{blue!10}{blue!55!black}{#1}{#2}{#3}}

% Espacio alrededor de tablas flotantes (mantiene el cuerpo en 3 páginas)
\setlength{\textfloatsep}{10pt plus 2pt minus 2pt}
\setlength{\intextsep}{8pt plus 2pt minus 2pt}
\captionsetup{skip=4pt}
\newcommand{\pp}{\,pp}
\newcolumntype{Y}{>{\raggedright\arraybackslash}X}

\begin{document}
\thispagestyle{primera}

% --- Título y datos ---
\noindent{\LARGE\bfseries Parada anticipada del MSAPSO\par}
\smallskip
\noindent{\large\textcolor{black!60}{Resumen ejecutivo · Más energía solar bajo sombra parcial con un ajuste de firmware}\par}
\medskip

\noindent\begingroup\setlength{\fboxsep}{7pt}%
\colorbox{black!8}{\parbox{\dimexpr\linewidth-2\fboxsep\relax}{\small\setstretch{1.1}%
\textcolor{black!60}{Integrantes:} P.\,C. Márquez Correo, D. Molina León, F. Moreira Guerra (Equipo 53 Solar Optimization)\\
\textcolor{black!60}{Programa:} Proyecto Integrador · Maestría en Inteligencia Artificial Aplicada\\
\textcolor{black!60}{Sponsor:} Prof. Gerardo J. Camacho González, PhD, y Tecnológico de Monterrey \quad \textcolor{black!60}{Fecha:} 14/11/2026}}%
\endgroup\par\medskip

% --- Problema y solución ---
Cuando parte de un arreglo fotovoltaico queda a la sombra, aparecen varios máximos de potencia y el controlador debe buscar el mayor midiendo en distintas corrientes \parencite{endiz2025}; mientras busca, el panel entrega menos de lo que podría. El algoritmo de referencia, el MSAPSO \parencite{hybridmppt2025}, combina un enjambre de partículas determinista \parencite{kennedy1995,rajendran2016} con un control Super-Twisting \parencite{levant2003}; es preciso, pero gasta siempre 175 mediciones por búsqueda y, con sombra cambiante, busca una y otra vez. Buscamos extraer más energía sin agregar sensores ni hardware. Comparamos con un mismo protocolo alternativas desde redes neuronales hasta optimización por entropía cruzada \parencite{rubinstein2004,tang2017,zhang2022}, y la mejor resultó ser la más sencilla: la \textbf{V7}, que conserva el MSAPSO y solo detiene la búsqueda en cuanto las partículas del enjambre se agrupan. No requiere entrenamiento ni datos: es un cambio de firmware.

% --- Cifras de desempeño del modelo ---
\medskip
\noindent
\begin{minipage}[t]{0.235\linewidth}\cifragris{+2.78\pp}{Energía entregada}{vs.\ MSAPSO; IC 95\,\%: 2.39–3.18}{}\end{minipage}\hfill
\begin{minipage}[t]{0.235\linewidth}\cifragris{$-$60\,\%}{Mediciones}{70 en lugar de 175 por búsqueda}{}\end{minipage}\hfill
\begin{minipage}[t]{0.235\linewidth}\cifragris{0.05\pp}{Pérdida de exactitud}{p99 de 0.64\pp\ (umbral: 3\pp)}{}\end{minipage}\hfill
\begin{minipage}[t]{0.235\linewidth}\cifragris{+6.7\pp}{Sombra frecuente}{ganancia en el régimen alto}{}\end{minipage}
\medskip

% --- Desempeño ---
Las cifras provienen de un simulador dinámico del convertidor y miden la energía entregada respecto del máximo alcanzable, en puntos porcentuales (pp), con intervalos de confianza por \textit{bootstrap} pareado \parencite{efron1993}. En 1\,200 escenarios de confirmación, nunca usados para ajustar parámetros, la V7 entrega el 87.38\,\% de la energía máxima contra el 84.60\,\% del MSAPSO, con 60\,\% menos mediciones y prácticamente la misma exactitud. La ventaja se mantuvo en todos los paneles, tiempos de espera y niveles de ruido evaluados. La Tabla~\ref{tab:metodos} muestra por qué no recomendamos una alternativa más compleja: el MLP + CEM, una red neuronal entrenada, gana 1.01\pp\ más en promedio, pero en el 1\,\% de sus peores casos pierde más de 11.7\pp\ de exactitud. La V14 gana 1.60\pp, pero su cola solo cumple con el límite de 12\,$\Omega$ y no en todos los paneles (p99 de 6.4\pp\ con 72 celdas): es una alternativa con condiciones.

\begin{table}[h]
\centering\footnotesize\setstretch{1}\setlength{\tabcolsep}{4pt}
\caption{Métodos en 1\,200 escenarios de confirmación. Energía en pp (IC 95\,\%). Cola p99 en 1\,000 curvas (umbral $\le 3$\pp; entre paréntesis, con límite de 12\,$\Omega$). Fuente: elaboración propia.}
\label{tab:metodos}
\begin{tabular}{@{}lccccc@{}}
\toprule
Método & Mediciones & Energía vs.\ MSAPSO & Energía vs.\ V7 & Pierde $>3$\pp & Cola p99\\
\midrule
MSAPSO (referencia) & 175 & --- & $-2.78$ [$-3.18$, $-2.39$] & --- & --- \\
\textbf{V7 (recomendada)} & \textbf{70} & $\mathbf{+2.78}$ [2.39, 3.18] & --- & 8\,\% & \textbf{0.64} (0.27) \\
Candidata V14 & 39 & $+4.38$ [3.95, 4.82] & $+1.60$ [1.24, 1.96] & 7\,\% & 13.0 (1.4) \\
MLP + CEM & 37 & $+3.80$ [3.19, 4.39] & $+1.01$ [0.49, 1.50] & 18\,\% & 11.7 (15.3) \\
DGBCO & 40 & $+3.64$ [3.13, 4.13] & $+0.85$ [0.47, 1.22] & 16\,\% & 15.9 (11.3) \\
DecentCEM & 40 & $+3.13$ [2.59, 3.65] & $+0.35$ [$-0.11$, 0.79] & 17\,\% & 5.3 (4.0) \\
CEDGO & 41 & $+1.98$ [1.39, 2.54] & $-0.81$ [$-1.32$, $-0.31$] & 22\,\% & 28.9 (26.9) \\
CEM puro & 59 & $+1.52$ [0.96, 2.09] & $-1.26$ [$-1.78$, $-0.77$] & 22\,\% & 20.1 (13.4) \\
\bottomrule
\end{tabular}
\end{table}

% --- Beneficios cuantificables (mensuales): mejora en el controlador ---
Siguiendo a \textcite{korolov2022}, medimos el beneficio en el desempeño del controlador. Con supuestos ilustrativos, no medidos en campo (6\,h de sol al día, 20\,\% de ese tiempo en sombra cambiante y una búsqueda por minuto en esos periodos, es decir, unas 2\,160 búsquedas al mes), los beneficios mensuales por controlador son (cálculo en el Anexo~A):

\medskip
\noindent
\begin{minipage}[t]{0.235\linewidth}\cifragris{0.22--0.70}{kWh/kWp al mes}{según la frecuencia de cambios de sombra}{}\end{minipage}\hfill
\begin{minipage}[t]{0.235\linewidth}\cifragris{$-$18.9\,min}{Ahorro de tiempo}{de búsqueda al mes}{}\end{minipage}\hfill
\begin{minipage}[t]{0.235\linewidth}\cifragris{+150\,\%}{Capacidad}{226\,800 mediciones/mes}{}\end{minipage}\hfill
\begin{minipage}[t]{0.235\linewidth}\cifragris{$-$56\,\%}{Menos errores}{8 vs.\ 18\,\% del MLP + CEM}{}\end{minipage}

\medskip

El controlador mejora sin cambiar el hardware: busca 2.5 veces más rápido, conserva la exactitud del MSAPSO y falla menos que las alternativas de aprendizaje (pérdida mayor de 3\pp\ en 8\,\% de las condiciones contra 18\,\%), y la mejora llega a cualquier controlador con enjambre mediante una actualización de firmware. Esperamos además beneficios intangibles: \textbf{mejora en la toma de decisiones}, porque el protocolo de evaluación (criterios previos, escenarios de confirmación y análisis del peor caso) permite decidir con evidencia sobre algoritmos futuros; \textbf{mejora en la calidad}, con exactitud al nivel del MSAPSO (pérdida máxima de 0.97\pp); \textbf{simplificación del trabajo}, sin modelo que entrenar y con una regla de unas 30 operaciones por pasada, e \textbf{innovación}, como mejora verificable sobre un algoritmo publicado y transferible a otros controladores.

% --- Inversión y costos (TCO a 12 meses) ---
Llevar la V7 a un equipo real requiere desarrollo (validarla contra la implementación del autor del MSAPSO y programarla en el controlador), implementación (prototipo, banco de pruebas \textit{hardware-in-the-loop} y piloto en un sitio con sombra cambiante) y seis meses de operación del piloto con telemetría en AWS IoT o Azure IoT Hub. El costo total de propiedad estimado a 12 meses, con 10\,\% de contingencia, es (desglose en el Anexo~B):

\medskip
\noindent
\begin{minipage}[t]{0.235\linewidth}\cifragris{\$90\,000}{Desarrollo}{MXN}{}\end{minipage}\hfill
\begin{minipage}[t]{0.235\linewidth}\cifragris{\$102\,000}{Implementación}{MXN}{}\end{minipage}\hfill
\begin{minipage}[t]{0.235\linewidth}\cifragris{\$54\,180}{Operación}{MXN, con contingencia}{}\end{minipage}\hfill
\begin{minipage}[t]{0.235\linewidth}\cifraacento{\$246\,180}{TCO 12 meses}{MXN}{}\end{minipage}
\medskip

\textbf{Sugerimos al asesor y al Tecnológico de Monterrey considerar la inversión en dos etapas.} La Etapa~1 (validación, firmware y banco de pruebas; 168\,000\,MXN) podría iniciarse en el corto plazo. La Etapa~2 (piloto, operación y contingencia; 78\,180\,MXN) convendría liberarla una vez que el banco de pruebas confirme tres condiciones: una ganancia de al menos 2\pp\ de energía en sombra cambiante, una cola con p99 $\le 3$\pp\ en cada panel, y una pérdida mayor de 3\pp\ en no más del 10\,\% de los escenarios. Una vez validada, cada controlador adicional cuesta casi nada, porque la mejora se distribuye como firmware.

% --- Riesgos y mitigación ---
\begin{table}[h]
\centering\footnotesize\setstretch{1}
\caption{Principales riesgos para avanzar y medidas de mitigación.}
\label{tab:riesgos}
\begin{tabularx}{\linewidth}{@{}p{0.47\linewidth}Y@{}}
\toprule
\textbf{Riesgo} & \textbf{Mitigación}\\
\midrule
Resultados solo de simulación. & La Etapa~1 valida contra la implementación del autor.\\
Sin ganancia en sombra estable ($-1.2$\pp). & Desplegar donde la sombra cambia con frecuencia.\\
En el 8\,\% de los escenarios pierde más de 3\pp. & Alertas sobre la cola; MSAPSO como respaldo.\\
Con ruido alto (2\,W) la ventaja baja a +1.59\pp. & Especificar sensores y alertar si el ruido excede lo validado.\\
El convertidor real difiere del simulado; p.\,ej., el enjambre arranca con 4--5 posiciones, no 7, porque las más bajas son inalcanzables. & Caracterizarlo en el banco y recalibrar; sembrar las partículas solo en la zona alcanzable.\\
\bottomrule
\end{tabularx}
\end{table}

% --- Transferencia y continuidad ---
Para dar continuidad, entregaremos al asesor y al Tecnológico de Monterrey el repositorio del proyecto (\url{https://github.com/pedro-mrqz/msapso-mppt}) con código y resultados, el simulador interactivo, el protocolo de evaluación, la especificación de la regla para firmware embebido de bajo costo \parencite{motahhir2020} y un plan de monitoreo con matriz de responsables \parencite{miller2022}, para que otro equipo continúe sin depender de los autores.

\section*{Declaración de uso de inteligencia artificial}
Se utilizó \textcite{anthropic2026claude} para implementar el simulador y analizar resultados, y \textcite{anthropic2026opus} para redactar este resumen. Los autores revisaron y asumen el contenido.

% --- Referencias (biblatex + Biber, solo lo citado) ---
\printbibliography[title={Referencias}]

% --- Anexos ---
\clearpage
\section*{Anexo A. Cálculo de los beneficios cuantificables}
Todos los beneficios se expresan por mes y por controlador (la productividad, por kWp). Son estimaciones ilustrativas a partir de los resultados de simulación; no se han medido en una instalación real.

\begin{table}[H]
\centering\small\setstretch{1.05}
\caption{Supuestos utilizados en la estimación de beneficios.}
\label{tab:supuestos}
\begin{tabularx}{\linewidth}{@{}Yrp{5.6cm}@{}}
\toprule
Supuesto & Valor & Fuente o justificación\\
\midrule
Ganancia de energía de la V7 en la confirmación & 2.78\pp & Resultados del proyecto, 1\,200 escenarios (IC 95\,\%: 2.39--3.18); corresponde a sombra que cambia como en los escenarios (del orden de 100 búsquedas por minuto); se usa para el extremo superior del rango\\
Producción anual específica & 1\,500 kWh/kWp & Supuesto ilustrativo del proyecto\\
Fracción del tiempo de operación en sombra cambiante & 20\,\% & Supuesto ilustrativo del proyecto; debe medirse en el sitio\\
Horas de sol al mes & 180 h & 6 h al día $\times$ 30 días; supuesto del equipo\\
Búsquedas por minuto en sombra cambiante & 1 & Supuesto del equipo; validar con el indicador de reinicios por hora\\
Mediciones por búsqueda (MSAPSO / V7) & 175 / 70 & Resultados del proyecto\\
Duración de una medición & 5 ms & Régimen de trabajo del simulador\\
Condiciones de sombra con pérdida mayor de 3\pp\ frente al MSAPSO (V7 / MLP + CEM) & 8\,\% / 18\,\% & Resultados del proyecto (Tabla~\ref{tab:metodos})\\
\bottomrule
\end{tabularx}
\end{table}

\begin{table}[H]
\centering\small\setstretch{1.05}
\caption{Cálculo de los indicadores mensuales.}
\label{tab:calculo}
\begin{tabularx}{\linewidth}{@{}>{\raggedright\arraybackslash}p{3.3cm}Y>{\raggedright\arraybackslash}p{3.2cm}@{}}
\toprule
Indicador & Base de cálculo & Valor mensual\\
\midrule
Búsquedas al mes (base) & $180\ \text{h} \times 0.20 \times 60 = 2\,160$ búsquedas & 2\,160 por controlador\\
Incremento en productividad & Con una búsqueda por minuto, la V7 ahorra $(175-70) \times 5\ \text{ms} = 0.525$ s por minuto, el 0.875\,\% del tiempo en sombra cambiante. Aun si el panel no entregara energía al buscar, la ganancia no supera 0.875\pp: $125 \times 0.20 \times 0.00875 = 0.219$, con $1\,500/12 = 125$ kWh/kWp al mes & Hasta +0.22 kWh/kWp con 1 búsqueda por minuto; $\approx$ 0.70 con la frecuencia de los escenarios\\
Ahorro de tiempo & MSAPSO: $2\,160 \times 175 \times 5\ \text{ms} = 31.5$ min; V7: $2\,160 \times 70 \times 5\ \text{ms} = 12.6$ min & $-$18.9 min por controlador ($-$60\,\%)\\
Capacidad de procesamiento & $175/70 = 2.5$ búsquedas en el tiempo de una; $2\,160 \times (175-70) = 226\,800$ mediciones liberadas & +150\,\%; 226\,800 mediciones\\
Reducción de errores & Fracción de condiciones de sombra con pérdida mayor de 3\pp: $(18-8)/18 = 0.56$; se mantiene cada mes mientras se repitan esas condiciones & $-$56\,\% frente al MLP + CEM\\
\bottomrule
\end{tabularx}
\end{table}

\noindent\textbf{Sensibilidad.} Los indicadores de tiempo y capacidad son proporcionales al número de búsquedas: con una búsqueda cada dos minutos, el ahorro de tiempo baja a 9.5 min al mes; con dos por minuto, sube a 37.8 min. Las proporciones ($-$60\,\%, +150\,\% y $-$56\,\%) no dependen de ese supuesto. La cota de productividad también es proporcional: 0.11 kWh/kWp al mes con una búsqueda cada dos minutos y 0.44 con dos por minuto; la ganancia real será menor, porque el panel sí entrega energía mientras busca. El extremo superior del rango corresponde a una sombra que cambia con la frecuencia de los escenarios evaluados (del orden de 100 búsquedas por minuto): ahí se alcanzan los 2.78\pp\ de la confirmación, unos 0.70 kWh/kWp al mes u 8 kWh/kWp al año, la cifra ilustrativa del artículo (0.60 al mes con el extremo inferior del intervalo, 2.39\pp). El beneficio real queda entre ambos extremos y se determinará en el piloto midiendo la frecuencia de cambios de sombra (reinicios por hora). Frente al MSAPSO, la V7 no reduce errores: su pérdida media de exactitud es de 0.05\pp, y en el 8\,\% de las condiciones pierde más de 3\pp. El ahorro de tiempo se refiere al tiempo que el controlador pasa buscando, no a horas-persona.

\section*{Anexo B. Costo total de propiedad (TCO) a 12 meses}
Supuestos: tarifa de 450 MXN por hora de ingeniería especializada; el laboratorio institucional proporciona el emulador de arreglo fotovoltaico a cambio de una cuota de uso; el piloto opera con hasta 10 equipos durante los meses 7 a 12. Las cifras son estimaciones del equipo y deben cotizarse antes de aprobar el gasto.

\begin{table}[H]
\centering\small\setstretch{1.05}
\caption{Desglose del TCO por servicio (MXN).}
\label{tab:tco}
\begin{tabularx}{\linewidth}{@{}lYp{3.4cm}r@{}}
\toprule
Etapa & Servicio & Base & Costo\\
\midrule
Desarrollo & Validación con la implementación del autor y datos medidos & 60 h $\times$ 450 & 27\,000\\
 & Programación de la V7 en el firmware del controlador & 100 h $\times$ 450 & 45\,000\\
 & Pruebas unitarias y verificación de tiempos de ejecución & 40 h $\times$ 450 & 18\,000\\
\addlinespace
Implementación & Prototipo: tarjeta de control, convertidor Boost y sensores & Compra única & 18\,000\\
 & Uso del emulador de arreglo FV (banco HIL) & 2 meses $\times$ 12\,000 & 24\,000\\
 & Ejecución de las pruebas HIL con los criterios de aceptación & 80 h $\times$ 450 & 36\,000\\
 & Instalación y puesta en marcha del piloto & 40 h $\times$ 450 + 6\,000 viáticos & 24\,000\\
\addlinespace
Operación & Telemetría y monitoreo en la nube (AWS IoT o Azure IoT Hub) & 6 meses $\times$ 500 & 3\,000\\
 & Almacenamiento de registros y tablero & 6 meses $\times$ 300 & 1\,800\\
 & Monitoreo, atención de alertas y mantenimiento & 6 meses $\times$ 10 h $\times$ 450 & 27\,000\\
\addlinespace
Contingencia & 10\,\% sobre el subtotal (223\,800) & & 22\,380\\
\midrule
\multicolumn{3}{@{}l}{\textbf{TCO 12 meses}} & \textbf{246\,180}\\
\multicolumn{3}{@{}l}{\quad Etapa 1: desarrollo, prototipo, banco HIL y pruebas} & 168\,000\\
\multicolumn{3}{@{}l}{\quad Etapa 2: piloto, operación y contingencia} & 78\,180\\
\bottomrule
\end{tabularx}
\end{table}

\noindent\textbf{Costo marginal.} Una vez validada la solución, el costo marginal por equipo adicional se limita a la distribución del firmware y a la telemetría (unos 50 MXN por equipo al mes).

\end{document}
